\documentclass[11pt,addpoints,answers]{exam}

%%%%%%%%%%%%%%%%%%%%%%%%%%%%%%%%%%%%%%%%%%%
% Commands for customizing the assignment %
%%%%%%%%%%%%%%%%%%%%%%%%%%%%%%%%%%%%%%%%%%%

\newcommand{\courseNum}{COURSE NUMBER}
\newcommand{\courseName}{COURSE NAME}
\newcommand{\courseSem}{COURSE SEMESTER}
\newcommand{\courseUrl}{\url{COURSE URL}}
\newcommand{\hwNum}{HOMEWORK NUMBER}
\newcommand{\hwTopic}{HOMEWORK TOPIC}
\newcommand{\hwName}{\hwNum: \hwTopic}
\newcommand{\outDate}{}
\newcommand{\dueDate}{}
\newcommand{\taNames}{TA NAMES}
\newcommand{\homeworktype}{\string written+prog}
\newcommand{\overleafUrl}{}

%% To HIDE SOLUTIONS (to post at the website for students), set this value to 0: \def\issoln{0}
\providecommand{\issoln}{0}
%\providecommand{\issoln}{1}

%-----------------------------------------------------------------------------
% PACKAGES AND OTHER DOCUMENT CONFIGURATIONS
%-----------------------------------------------------------------------------

\usepackage[margin=1in]{geometry}
\usepackage{amsmath, amsfonts}
\usepackage{enumerate}
\usepackage{graphicx}
\usepackage{titling}
\usepackage{url}
\usepackage{xfrac}
\usepackage{natbib}
\usepackage{amssymb}
\usepackage{amsthm}
\usepackage{paralist}
\usepackage{epstopdf}
\usepackage{tabularx}
\usepackage{longtable}
\usepackage{multirow}
\usepackage{multicol}
\usepackage[colorlinks=true,urlcolor=blue]{hyperref}
\usepackage{algorithm}
\usepackage{algorithmicx}
\usepackage[noend]{algpseudocode}
\usepackage{float}
\usepackage{enumerate}
\usepackage{array}
\usepackage{environ}
\usepackage{times}
\usepackage{textcomp}
\usepackage{caption}
\usepackage{parskip} % For NIPS style paragraphs.
\usepackage[compact]{titlesec} % Less whitespace around titles
\usepackage[inline]{enumitem} % For inline enumerate* and itemize*
\usepackage{datetime}
\usepackage{comment}
% \usepackage{minted}
\usepackage{lastpage}
\usepackage{color}
% \usepackage{xcolor}
\usepackage[dvipsnames]{xcolor}
\usepackage[final]{listings}
\usepackage{tikz}
\usetikzlibrary{shapes,decorations}
\usepackage{framed}
\usepackage{booktabs}
\usepackage{cprotect}
\usepackage{verbatimbox}
\usepackage{multicol}
\usepackage{hyperref}
\usepackage{subcaption}
\usepackage{mathtools} % For drcases
\usepackage{cancel}
\usepackage[many]{tcolorbox}
\usepackage{soul}
\usepackage[bottom]{footmisc}
\usepackage{bm}
\usepackage{wasysym}
\usepackage{lipsum}

\usepackage{tikz}
\usetikzlibrary{arrows}
\usetikzlibrary{arrows.meta}
\usetikzlibrary{shapes.geometric}
\usetikzlibrary{positioning, arrows, automata, calc}
\usepackage{transparent}
\usepackage{tikz-cd}

%%%%%%%%%%%%%%%%%%%%%%%%%%%%%%%%%%%%%%%%%%%
% Formatting for \CorrectChoice of "exam" %
%%%%%%%%%%%%%%%%%%%%%%%%%%%%%%%%%%%%%%%%%%%

\CorrectChoiceEmphasis{}
\checkedchar{\blackcircle}

\newenvironment{checkboxessquare}{
    \begingroup
    \checkboxchar{$\Box$} \checkedchar{$\blacksquare$} % change checkbox style locally
    \begin{checkboxes}
    }{
    \end{checkboxes}
    \endgroup
    }

%%%%%%%%%%%%%%%%%%%%%%%%%%%%%%%%%%%%%%%%%%%
% Better numbering                        %
%%%%%%%%%%%%%%%%%%%%%%%%%%%%%%%%%%%%%%%%%%%

% \numberwithin{equation}{section} % Number equations within sections (i.e. 1.1, 1.2, 2.1, 2.2 instead of 1, 2, 3, 4)
% \numberwithin{figure}{section} % Number figures within sections (i.e. 1.1, 1.2, 2.1, 2.2 instead of 1, 2, 3, 4)
% \numberwithin{table}{section} % Number tables within sections (i.e. 1.1, 1.2, 2.1, 2.2 instead of 1, 2, 3, 4)

%%%%%%%%%%%%%%%%%%%%%%%%%%%%%%%%%%%%%%%%%%
% Custom commands                        %
%%%%%%%%%%%%%%%%%%%%%%%%%%%%%%%%%%%%%%%%%%
\newcommand{\blackcircle}{\tikz\draw[black,fill=black] (0,0) circle (1ex);}
\renewcommand{\circle}{\tikz\draw[black] (0,0) circle (1ex);}

\newcommand{\solo}{ \textcolor{orange}{[SOLO]} }
\newcommand{\open}{ \textcolor{blue}{[OPEN]} }


%%%%%%%%%%%%%%%%%%%%%%%%%%%%%%%%%%%%%%%%%%
% Custom commands for Math               %
%%%%%%%%%%%%%%%%%%%%%%%%%%%%%%%%%%%%%%%%%%
\newcommand{\vc}[1]{\boldsymbol{#1}}
\newcommand{\adj}[1]{\frac{\partial \ell}{\partial #1}}
\newcommand{\chain}[2]{\adj{#2} = \adj{#1}\frac{\partial #1}{\partial #2}}
\newcommand{\ntset}{test}
\newcommand{\zerov}{\mathbf{0}}
\DeclareMathOperator*{\argmin}{argmin}

% mathcal
\newcommand{\Ac}{\mathcal{A}}
\newcommand{\Bc}{\mathcal{B}}
\newcommand{\Cc}{\mathcal{C}}
\newcommand{\Dc}{\mathcal{D}}
\newcommand{\Ec}{\mathcal{E}}
\newcommand{\Fc}{\mathcal{F}}
\newcommand{\Gc}{\mathcal{G}}
\newcommand{\Hc}{\mathcal{H}}
\newcommand{\Ic}{\mathcal{I}}
\newcommand{\Jc}{\mathcal{J}}
\newcommand{\Kc}{\mathcal{K}}
\newcommand{\Lc}{\mathcal{L}}
\newcommand{\Mc}{\mathcal{M}}
\newcommand{\Nc}{\mathcal{N}}
\newcommand{\Oc}{\mathcal{O}}
\newcommand{\Pc}{\mathcal{P}}
\newcommand{\Qc}{\mathcal{Q}}
\newcommand{\Rc}{\mathcal{R}}
\newcommand{\Sc}{\mathcal{S}}
\newcommand{\Tc}{\mathcal{T}}
\newcommand{\Uc}{\mathcal{U}}
\newcommand{\Vc}{\mathcal{V}}
\newcommand{\Wc}{\mathcal{W}}
\newcommand{\Xc}{\mathcal{X}}
\newcommand{\Yc}{\mathcal{Y}}
\newcommand{\Zc}{\mathcal{Z}}

% mathbb
\newcommand{\Ab}{\mathbb{A}}
\newcommand{\Bb}{\mathbb{B}}
\newcommand{\Cb}{\mathbb{C}}
\newcommand{\Db}{\mathbb{D}}
\newcommand{\Eb}{\mathbb{E}}
\newcommand{\Fb}{\mathbb{F}}
\newcommand{\Gb}{\mathbb{G}}
\newcommand{\Hb}{\mathbb{H}}
\newcommand{\Ib}{\mathbb{I}}
\newcommand{\Jb}{\mathbb{J}}
\newcommand{\Kb}{\mathbb{K}}
\newcommand{\Lb}{\mathbb{L}}
\newcommand{\Mb}{\mathbb{M}}
\newcommand{\Nb}{\mathbb{N}}
\newcommand{\Ob}{\mathbb{O}}
\newcommand{\Pb}{\mathbb{P}}
\newcommand{\Qb}{\mathbb{Q}}
\newcommand{\Rb}{\mathbb{R}}
\newcommand{\Sb}{\mathbb{S}}
\newcommand{\Tb}{\mathbb{T}}
\newcommand{\Ub}{\mathbb{U}}
\newcommand{\Vb}{\mathbb{V}}
\newcommand{\Wb}{\mathbb{W}}
\newcommand{\Xb}{\mathbb{X}}
\newcommand{\Yb}{\mathbb{Y}}
\newcommand{\Zb}{\mathbb{Z}}

% mathbf lowercase
\newcommand{\av}{\mathbf{a}}
\newcommand{\bv}{\mathbf{b}}
\newcommand{\cv}{\mathbf{c}}
\newcommand{\dv}{\mathbf{d}}
\newcommand{\ev}{\mathbf{e}}
\newcommand{\fv}{\mathbf{f}}
\newcommand{\gv}{\mathbf{g}}
\newcommand{\hv}{\mathbf{h}}
\newcommand{\iv}{\mathbf{i}}
\newcommand{\jv}{\mathbf{j}}
\newcommand{\kv}{\mathbf{k}}
\newcommand{\lv}{\mathbf{l}}
\newcommand{\mv}{\mathbf{m}}
\newcommand{\nv}{\mathbf{n}}
\newcommand{\ov}{\mathbf{o}}
\newcommand{\pv}{\mathbf{p}}
\newcommand{\qv}{\mathbf{q}}
\newcommand{\rv}{\mathbf{r}}
\newcommand{\sv}{\mathbf{s}}
\newcommand{\tv}{\mathbf{t}}
\newcommand{\uv}{\mathbf{u}}
\newcommand{\vv}{\mathbf{v}}
\newcommand{\wv}{\mathbf{w}}
\newcommand{\xv}{\mathbf{x}}
\newcommand{\yv}{\mathbf{y}}
\newcommand{\zv}{\mathbf{z}}

% mathbf uppercase
\newcommand{\Av}{\mathbf{A}}
\newcommand{\Bv}{\mathbf{B}}
\newcommand{\Cv}{\mathbf{C}}
\newcommand{\Dv}{\mathbf{D}}
\newcommand{\Ev}{\mathbf{E}}
\newcommand{\Fv}{\mathbf{F}}
\newcommand{\Gv}{\mathbf{G}}
\newcommand{\Hv}{\mathbf{H}}
\newcommand{\Iv}{\mathbf{I}}
\newcommand{\Jv}{\mathbf{J}}
\newcommand{\Kv}{\mathbf{K}}
\newcommand{\Lv}{\mathbf{L}}
\newcommand{\Mv}{\mathbf{M}}
\newcommand{\Nv}{\mathbf{N}}
\newcommand{\Ov}{\mathbf{O}}
\newcommand{\Pv}{\mathbf{P}}
\newcommand{\Qv}{\mathbf{Q}}
\newcommand{\Rv}{\mathbf{R}}
\newcommand{\Sv}{\mathbf{S}}
\newcommand{\Tv}{\mathbf{T}}
\newcommand{\Uv}{\mathbf{U}}
\newcommand{\Vv}{\mathbf{V}}
\newcommand{\Wv}{\mathbf{W}}
\newcommand{\Xv}{\mathbf{X}}
\newcommand{\Yv}{\mathbf{Y}}
\newcommand{\Zv}{\mathbf{Z}}

% bold greek lowercase
\newcommand{\alphav     }{\boldsymbol \alpha     }
\newcommand{\betav      }{\boldsymbol \beta      }
\newcommand{\gammav     }{\boldsymbol \gamma     }
\newcommand{\deltav     }{\boldsymbol \delta     }
\newcommand{\epsilonv   }{\boldsymbol \epsilon   }
\newcommand{\varepsilonv}{\boldsymbol \varepsilon}
\newcommand{\zetav      }{\boldsymbol \zeta      }
\newcommand{\etav       }{\boldsymbol \eta       }
\newcommand{\thetav     }{\boldsymbol \theta     }
\newcommand{\varthetav  }{\boldsymbol \vartheta  }
\newcommand{\iotav      }{\boldsymbol \iota      }
\newcommand{\kappav     }{\boldsymbol \kappa     }
\newcommand{\varkappav  }{\boldsymbol \varkappa  }
\newcommand{\lambdav    }{\boldsymbol \lambda    }
\newcommand{\muv        }{\boldsymbol \mu        }
\newcommand{\nuv        }{\boldsymbol \nu        }
\newcommand{\xiv        }{\boldsymbol \xi        }
\newcommand{\omicronv   }{\boldsymbol \omicron   }
\newcommand{\piv        }{\boldsymbol \pi        }
\newcommand{\varpiv     }{\boldsymbol \varpi     }
\newcommand{\rhov       }{\boldsymbol \rho       }
\newcommand{\varrhov    }{\boldsymbol \varrho    }
\newcommand{\sigmav     }{\boldsymbol \sigma     }
\newcommand{\varsigmav  }{\boldsymbol \varsigma  }
\newcommand{\tauv       }{\boldsymbol \tau       }
\newcommand{\upsilonv   }{\boldsymbol \upsilon   }
\newcommand{\phiv       }{\boldsymbol \phi       }
\newcommand{\varphiv    }{\boldsymbol \varphi    }
\newcommand{\chiv       }{\boldsymbol \chi       }
\newcommand{\psiv       }{\boldsymbol \psi       }
\newcommand{\omegav     }{\boldsymbol \omega     }

% bold greek uppercase
\newcommand{\Gammav     }{\boldsymbol \Gamma     }
\newcommand{\Deltav     }{\boldsymbol \Delta     }
\newcommand{\Thetav     }{\boldsymbol \Theta     }
\newcommand{\Lambdav    }{\boldsymbol \Lambda    }
\newcommand{\Xiv        }{\boldsymbol \Xi        }
\newcommand{\Piv        }{\boldsymbol \Pi        }
\newcommand{\Sigmav     }{\boldsymbol \Sigma     }
\newcommand{\Upsilonv   }{\boldsymbol \Upsilon   }
\newcommand{\Phiv       }{\boldsymbol \Phi       }
\newcommand{\Psiv       }{\boldsymbol \Psi       }
\newcommand{\Omegav     }{\boldsymbol \Omega     }

%%%%%%%%%%%%%%%%%%%%%%%%%%%%%%%%%%%%%%%%%%%
% Code highlighting with listings         %
%%%%%%%%%%%%%%%%%%%%%%%%%%%%%%%%%%%%%%%%%%%

\definecolor{bluekeywords}{rgb}{0.13,0.13,1}
\definecolor{greencomments}{rgb}{0,0.5,0}
\definecolor{redstrings}{rgb}{0.9,0,0}
\definecolor{light-gray}{gray}{0.95}

\newcommand{\MYhref}[3][blue]{\href{#2}{\color{#1}{#3}}}%

\definecolor{dkgreen}{rgb}{0,0.6,0}
\definecolor{gray}{rgb}{0.5,0.5,0.5}
\definecolor{mauve}{rgb}{0.58,0,0.82}

\lstdefinelanguage{Shell}{
  keywords={tar, cd, make},
  %keywordstyle=\color{bluekeywords}\bfseries,
  alsoletter={+},
  ndkeywords={python3, python, py, javac, java, gcc, c, g++, cpp, .txt, octave, m, .tar},
  %ndkeywordstyle=\color{bluekeywords}\bfseries,
  identifierstyle=\color{black},
  sensitive=false,
  comment=[l]{//},
  morecomment=[s]{/*}{*/},
  commentstyle=\color{purple}\ttfamily,
  %stringstyle=\color{red}\ttfamily,
  morestring=[b]',
  morestring=[b]",
  backgroundcolor = \color{light-gray}
}

\lstset{columns=fixed, basicstyle=\ttfamily,
    backgroundcolor=\color{light-gray},xleftmargin=0.5cm,frame=tlbr,framesep=4pt,framerule=0pt}

\newcommand{\emptysquare}{{\LARGE $\square$}\ \ }
\newcommand{\filledsquare}{{\LARGE $\boxtimes$}\ \ }
\def \ifempty#1{\def\temp{#1} \ifx\temp\empty }


% \newcommand{\squaresolutionspace}[2][\emptysquare]{\newline #1}{#2}
\def \squaresolutionspace#1{ \ifempty{#1} \emptysquare \else #1\hspace{0.75pt}\fi}


\newcommand{\emptycircle}{{\LARGE $\fullmoon$}\ \ }
\newcommand{\filledcircle}{{\LARGE $\newmoon$}\ \ }
\def \circlesolutionspace#1{ \ifempty{#1} \emptycircle \else #1\hspace{0.75pt}\fi}
%%%%%%%%%%%%%%%%%%%%%%%%%%%%%%%%%%%%%%%%%%%
% Custom box for highlights               %
%%%%%%%%%%%%%%%%%%%%%%%%%%%%%%%%%%%%%%%%%%%

% Define box and box title style
\tikzstyle{mybox} = [fill=blue!10, very thick,
    rectangle, rounded corners, inner sep=1em, inner ysep=1em]

\NewEnviron{notebox}{
\begin{tikzpicture}
\node [mybox] (box){
    \begin{minipage}{0.95\textwidth}
        \BODY
    \end{minipage}
};
\end{tikzpicture}
}

%%%%%%%%%%%%%%%%%%%%%%%%%%%%%%%%%%%%%%%%%%%
% Commands showing / hiding solutions     %
%%%%%%%%%%%%%%%%%%%%%%%%%%%%%%%%%%%%%%%%%%%
\newcommand{\solutionspace}[4]{\fbox{\begin{minipage}[t][#1][t]{#2} \textbf{#3} \solution{}{#4} \end{minipage}}}

% To HIDE SOLUTIONS, set this value to 0:
%\providecommand{\issoln}{0}
%\providecommand{\issoln}{1}

\ifthenelse{\equal{\issoln}{1}}{

% SOLUTION environment
\newenvironment{soln}{\leavevmode\color{red}\ignorespaces }{}

% QUESTION AUTHORS environment
\newenvironment{qauthor}{\leavevmode\color{blue}\ignorespaces }{}

% Question tester comment environment
\newenvironment{qtester}{\leavevmode\color{green}\ignorespaces}{}

% Question learning objective comment environment
\newenvironment{qlearningobjective}{\leavevmode\color{green}\ignorespaces}{}

}{ % ELSE

  \NewEnviron{soln}{}
  \NewEnviron{qauthor}{}
  \NewEnviron{qtester}{}
  \NewEnviron{qlearningobjective}{}

}

%% To HIDE TAGS set this value to 0:
\def\showtags{0}
%%%%%%%%%%%%%%%%
\ifcsname showtags\endcsname \else \def\showtags{1} \fi
% Default to an empty tags environ.
\NewEnviron{tags}{}{}
\if\showtags 1
% Otherwise, include solutions as below.
\RenewEnviron{tags}{
    \fbox{
    \leavevmode\color{blue}\ignorespaces
    \textbf{TAGS:} \texttt{\url{\BODY}}
    }
    \vspace{-.5em}
}{}
\fi

\newtcolorbox[]{answer_box}[1][]{
    % breakable,
    fit,
    enhanced,
    % nobeforeafter,
    colback=white,
    title=Your Answer,
    sidebyside align=top,
    box align=top,
    #1
}

%\pagestyle{fancyplain}
\lhead{\hwName{}
}
\rhead{\courseNum}
\cfoot{\thepage{} of \numpages{}}

\title{\textsc{\hwNum \footnote{Adapted from CMU 10-623 Generative AI Homework Template}}
\\
\textsc{\hwTopic}
\thanks{Compiled on \today{} at \currenttime{}}\\
\vspace{1em}
} % Title

\author{\textsc{\large \courseNum{} \courseName{}}\\
\courseUrl
\vspace{1em}\\
\ifdefempty{\outDate}{}{  OUT: \outDate \\ }
\ifdefempty{\dueDate}{}{  DUE: \dueDate \\ }
  TAs: \taNames{}
}

\date{}

%%%%%%%%%%%%%%%%%%%%%%%%%%%%%%%%%%%%%%%%%%%%%%%%%
% Useful commands for typesetting the questions %
%%%%%%%%%%%%%%%%%%%%%%%%%%%%%%%%%%%%%%%%%%%%%%%%%

% This command will allow long \lstinline{} text to wrap automatically.
\sloppy

%%%%%%%%%%%%%%%%%%%%%%%%%%
% Document configuration %
%%%%%%%%%%%%%%%%%%%%%%%%%%

% Don't display a date in the title and remove the white space
\predate{}
\postdate{}
\date{}

% Don't display an author and remove the white space
% \preauthor{}
% \postauthor{}

% Question type commands
\newcommand{\sall}{\textbf{Select all that apply: }}
\newcommand{\sone}{\textbf{Select one: }}
\newcommand{\tf}{\textbf{True or False: }}




% Changes to examdoc
\qformat{\textbf{{\Large \thequestion \; \; \thequestiontitle \ (\totalpoints \ points)}} \hfill}
\renewcommand{\thequestion}{\arabic{question}}
\renewcommand{\questionlabel}{\thequestion.}

\renewcommand{\thepartno}{\thequestion.\arabic{partno}}
%\renewcommand{\partlabel}{\thequestion.\thepartno.}
\renewcommand{\partlabel}{\thepartno.}

% not working: \renewcommand{\subpartlabel}{(\thequestion.\thepartno.\thesubpart)}
% Commented after adding \question.\thepartno.
%\renewcommand{\partshook}{\setlength{\leftmargin}{0pt}}

\renewcommand{\thesubpart}{\thepartno.\alph{subpart}}
\renewcommand{\subpartlabel}{\thesubpart.}

\renewcommand{\thesubsubpart}{\thesubpart.\roman{subsubpart}}
\renewcommand{\subsubpartlabel}{\thesubsubpart.}

% copied from stack overflow, as all good things are
\newcommand\invisiblesection[1]{%
  \refstepcounter{section}%
  \addcontentsline{toc}{section}{\protect\numberline{\thesection}#1}%
  \sectionmark{#1}}

% quite possibly the worst workaround i have made for this class
\newcommand{\sectionquestion}[1]{
\titledquestion{#1}
\invisiblesection{#1}
~\vspace{-1em}
}

%%%%%%%%%%%%%%%%%%%%%%%%%%%%%%%%%%%%%%%%%%%
% New Environment for Pseudocode          %
%%%%%%%%%%%%%%%%%%%%%%%%%%%%%%%%%%%%%%%%%%%

% Python style for highlighting
\DeclareFixedFont{\ttb}{T1}{txtt}{bx}{n}{12} % for bold
\DeclareFixedFont{\ttm}{T1}{txtt}{m}{n}{12}  % for normal

\definecolor{deepblue}{rgb}{0,0,0.5}
\definecolor{deepred}{rgb}{0.6,0,0}
\definecolor{deepgreen}{rgb}{0,0.5,0}

\newcommand\pythonstyle{\lstset{
language=Python,
basicstyle=\ttm,
morekeywords={self},              % Add keywords here
keywordstyle=\ttb\color{deepblue},
emph={MyClass,__init__},          % Custom highlighting
emphstyle=\ttb\color{deepred},    % Custom highlighting style
stringstyle=\color{deepgreen},
frame=tb,                         % Any extra options here
showstringspaces=false
}}


% Python environment
\lstnewenvironment{your_code_solution}[1][]
{
\pythonstyle
\lstset{#1}
}
{}
\newcommand{\revert}[1]{\textcolor{orange}{#1}}

\begin{document}
 
\maketitle 

\newcommand \maxsubs {10 }
\section*{Instructions}
\begin{itemize}

\item \textbf{Collaboration Policy}: Please read the collaboration policy in the syllabus.
\item\textbf{Late Submission Policy:} See the late submission policy in the syllabus.
\item\textbf{Submitting your work:} You will use Gradescope to submit
  answers to all questions\ifthenelse{\equal{\homeworktype}{\string written}}{}{ and code}. 

\begin{itemize}
    
    \item \textbf{Written:} You will submit your completed homework as a PDF to
    Gradescope. Please use the provided template for the submissions which can
    be written in \LaTeX{} (via \href{www.overleaf.com}{Overleaf}) or directly
    on the PDF document. Each answer should be within the box provided. If you
    do not follow the template, your assignment may not be graded correctly and
    there will be a \textbf{\textcolor{red}{2\% penalty}} (e.g., if the homework
    is out of 100 points, 2 points will be deducted from your final score).
    
    \ifdefempty{\overleafUrl}{}{
    \item \textbf{\LaTeX{} Source:} \overleafUrl
    }

    \ifthenelse{\equal{\homeworktype}{\string written}}{}{
    % \item \textbf{Programming:} You will submit your code for programming
    % questions to Gradescope. There is no autograder. We will examine your code
    % by hand and may award marks for its submission.
    % }{}
    \item \textbf{Programming:} INSTRUCTIONS FOR PROGRAMMING SUBMISSION
    }{}
   
  \end{itemize}
  
\ifthenelse{\equal{\homeworktype}{\string written}}{}{\item\textbf{Materials:} The data that you will need in order to complete this assignment is posted along with the writeup and template on the course website.}

\end{itemize}

\begin{center}
    \pointtable[v][questions]
\end{center}\clearpage

\section*{Introduction}
\label{sec:intro}
Introduction section for homework

At a high-level, you will proceed as follows:

\begin{enumerate}
\item TASK 1
\item TASK 2
\item TASK 3
\end{enumerate}

\clearpage

\section*{Computing Environment}
\label{sec:env}

First you need to setup your computing environment, this is useful for future
homework assignments. Below we outline how you could do so on your machine.

\subsection*{Local Environment}
\label{sec:local}

For this course, we suggest using the \texttt{uv} project and package manager as
it provides a lightweight solution to automatically handle many of the
dependency and python versioning challenges encountered with setups such as
Anaconda, MiniConda, etc. If you're familiar with tools such as \texttt{pip} and
\texttt{venv}, then \texttt{uv} will be pretty similar to what you're used to.

You can explore all of the available commands using \texttt{uv --help} but in
short, here are a couple of useful commands that you'll often use:

\begin{table}[H]
    \centering
    \begin{tabular}{|l|l|}
        \hline
        \texttt{source .venv/bin/activate} & (Linux and Mac) Activates python environment. \\
        \hline
        \texttt{source .venv/Scripts/Activate} & (Windows) Activates python environment. \\
        \hline
        \texttt{uv sync} & Installs  / synchronizes projects dependencies. \\
        \hline
        \texttt{uv add <package>} & Installs \texttt{pip} package and adds it to project dependencies. \\
        \hline
        \texttt{uv remove <package>} & Removes \texttt{pip} package dependency from project. \\
        \hline
        \texttt{uv run <script.py>} & Runs python script inside of \texttt{uv} environment. \\
        \hline
    \end{tabular}
\end{table}

\begin{enumerate}
\item Follow the instructions linked below to install UV.  

\url{https://docs.astral.sh/uv/}

\item Navigate to folder \texttt{homework-0-course-primer/} and synchronize dependencies with:
\begin{verbatim}
uv sync
\end{verbatim}

\item Activate the environment

Mac and Linux:
\begin{verbatim}
source .venv/bin/activate
\end{verbatim}

Windows:
\begin{verbatim}
source .venv/Scripts/activate
\end{verbatim}

\item Test \text{uv} works and run \texttt{main.py} file.
\begin{verbatim}
uv run main.py
\end{verbatim}

\end{enumerate}

\begin{questions}

\clearpage

\sectionquestion{Getting Started}

\begin{parts}
\part[5] Did you read the syllabus?

    \begin{checkboxes}
     % \correctchoice Yes 
     \choice Yes 
     \choice No
    \end{checkboxes}
    
\part[5] Create or provide the username to your HuggingFace account 

    \begin{answer_box}[title=,height=1cm, width=10cm]

    \end{answer_box}

\part[5] Did you successfully install \texttt{uv} onto your machine?

    \begin{checkboxes}
     % \correctchoice Yes 
     \choice Yes 
     \choice No
    \end{checkboxes}
    
\end{parts}

\clearpage

\sectionquestion{Example LLM Question}

In this section, you will create a Large Language Model (LLM) from ``scratch"
following the instruction set from HuggingFace's \textit{Training a Causal
Language Model from Scratch} tutorial.

\url{https://huggingface.co/learn/llm-course/en/chapter7/6}

We will follow most of the steps provided in this tutorial as it guides us on
how to use HuggingFace's abstractions for dataset and model management. In
addition, a couple more tasks will help refresh us on some of the basic concepts
of how an LLM is trained and inferenced upon. Since most of the code is already
provided to us, we'll look deeper into what's going on under the hood.

\subsection*{Dataset}

Let's take a look at the \texttt{codeparrot} dataset hosted on HuggingFace

\url{https://huggingface.co/datasets/transformersbook/codeparrot}

\begin{parts}

\part[5] Briefly describe this dataset. What kind of data is in this dataset
where does the majority of its content come from?

    \begin{answer_box}[title=,height=3cm, width=15cm]


    \end{answer_box}
    
\part It seems that this dataset consists of a large amount of data.

\begin{subparts}
    
    \subpart[1] What do the authors of this dataset claim its uncompressed size to be? 
        
        \begin{answer_box}[title=,height=1cm, width=5cm]
        
        \end{answer_box}

    \subpart[1] How was this dataset obtained?
    
    \textit{Let's go a bit deeper than just saying Google BigQuery}
        \begin{answer_box}[title=,height=3cm, width=15cm]
    
        \end{answer_box}

\end{subparts}

\part Now, lets look at the paper \textit{Scaling Laws for Neural Language
Models} by \cite{kaplan_scaling_2020}

\begin{subparts}

    \subpart[5] What relationship do the authors observe between model size and
    language modeling performance, and how is this relationship characterized
    mathematically?
    
        \begin{answer_box}[title=,height=3cm, width=15cm]
        
        \end{answer_box}

\end{subparts}

\subsection*{Tokenization}
For this section we'll look at the tokenization for processing the dataset along
with some concepts from the reading \textit{Attention is all you need}
\cite{vaswani_attention_2023}.

\part[5] Briefly describe how the dataset is tokenized and prepared for training
in this portion of the tutorial. Why split documents into fixed-length chunks
instead of truncating them?

    \begin{answer_box}[title=,height=3cm, width=15cm]

    \end{answer_box}
    
\part The tutorial fixes a relatively small context size and applies chunking
during tokenization.

\begin{subparts}
    
    \subpart[2] What context length (in tokens) is used for training in this
    tutorial?
        
        \begin{answer_box}[title=,height=1cm, width=5cm]
        
        \end{answer_box}
        
\end{subparts}

\part[5] The tutorial notes that increasing the context length improves the
amount of information available to the model but increases computational cost.
Explain how the choice of context length affects both training efficiency and
memory usage, and relate this tradeoff to the self-attention mechanism
introduced in \textit{Attention Is All You Need} \cite{vaswani_attention_2023}.

    \begin{answer_box}[title=,height=3cm, width=15cm]

    \end{answer_box}

\end{parts}

\clearpage

\sectionquestion{Code Upload}

\begin{parts}

\part[0] Did you upload your code to the appropriate programming slot on Gradescope? \\
\emph{Hint:} The correct answer is `yes'.
    \begin{checkboxes}
        % \correctchoice Yes 
        \choice Yes 
        \choice No
    \end{checkboxes}

\end{parts}

\newpage

\sectionquestion{Collaboration and Generative AI Questions}

\begin{parts}

\uplevel{
After you have completed all other components of this assignment, report your
answers to these questions regarding the collaboration policy. Details of the
policy can be found in the syllabus.
}

    \part[2] Did you collaborate with anyone on this assignment? If so, list
    their name or Andrew ID and which problems you worked together on.

        \begin{answer_box}[title=,height=3cm, width=15cm]
        \end{answer_box}

    \part[3] Did you use generative AI for parts of this assignment? If so,
    include the input prompts along with links to the conversation.
    
        \begin{answer_box}[title=,height=3cm, width=15cm]
        \end{answer_box}
        
\end{parts}
\end{questions}

\bibliographystyle{plainnat}
\bibliography{references}

\end{document}
